\documentclass[11pt,a4paper]{article}
\usepackage[T1]{fontenc}
\usepackage{lmodern,amsmath,amssymb,amsthm,mathtools,microtype}
\usepackage[margin=26mm]{geometry}
\usepackage[colorlinks=true,linkcolor=blue,citecolor=blue,urlcolor=blue]{hyperref}
\newtheorem{theorem}{Theorem}[section]
\newtheorem{proposition}[theorem]{Proposition}
\newtheorem{lemma}[theorem]{Lemma}
\newtheorem{corollary}[theorem]{Corollary}
\theoremstyle{definition}
\newtheorem{definition}[theorem]{Definition}
\newtheorem{remark}[theorem]{Remark}
\newcommand{\Q}{\mathbb Q}
\newcommand{\Z}{\mathbb Z}
\newcommand{\hF}{h_{\mathrm F}}
\newcommand{\rad}{\operatorname{rad}}
\newcommand{\ord}{\operatorname{ord}}
\newcommand{\Jac}{\operatorname{Jac}}
\setlength{\emergencystretch}{2em}
\title{An Explicit Five-Contact Height Bound of Exponent Eighteen}
\author{}
\date{9 November 2026}
\begin{document}
\maketitle
\begin{abstract}
We derive an explicit bound $h(a)\le H_{18}M_5(a)^{18}$ for the
logarithmic height of a rational parameter in terms of its odd
five-contact radical. The arithmetic input is the modular realization
used in the Diophantine construction. Retaining the Euler totient in
the order of its scalar character gives a cubic bound for the degree
of the relevant modular curve. A general Belyi-curve height estimate
and an explicit isogeny bound then give exponent eighteen, with
$H_{18}<2^{1984521}$. No additional $abc$ or Szpiro hypothesis is used.
We also give the corresponding exponent $18/r$ for $r$ disjoint
sets of contacts. All variables, exceptional primes, and numerical
constants are specified.
\end{abstract}

\section{Heights and the arithmetic input}

For $a=u/v\in\Q$ in lowest terms, with $v>0$, its absolute
logarithmic Weil height is
\[
 h(a)=\log\max(|u|,v).
\]
For an abelian variety $A$ over a number field, $\hF(A)$ denotes
its stable Faltings height in the original Faltings normalization.
Put
\[
 S_*=\{p\text{ prime}:p\le76176\},\qquad
 B_*=\prod_{p\in S_*}p,\qquad
 \mathcal D=\{0,1,2401,-343/32\},\qquad C=210^{300}.
\]
For $a=u/v\in\Q\setminus\mathcal D$, let $M_5(a)$ be the
product of the primes $p\notin S_*$ for which at least one of
\begin{equation}\label{eq:contacts}
 v_p(v),\quad v_p(u),\quad v_p(u-v),\quad
 v_p(u-2401v),\quad v_p(32u+343v)
\end{equation}
is positive and odd. The valuation satisfies $v_p(p)=1$.
The first form is the contact at infinity. Each prime occurs only
once, so $M_5(a)$ is a positive squarefree integer.

We use the following precise part of the arithmetic construction in
\cite[proof of \texttt{o8-height}, equations
\textup{(LR2)--(LR5)}, and proof of
\texttt{quartic-height19}]{Glazer}.
Stating this input explicitly separates the height calculation from
the geometric and modular realization on which it depends.

\begin{definition}[Modular realization input]\label{def:input}
For each $a\in\Q\setminus\mathcal D$, set $R=B_*M_5(a)$.
There is an abelian surface $A$ such that
\begin{equation}\label{eq:geometric}
 h(a)\le2\max(0,\hF(A))+128.
\end{equation}
Either $\hF(A)<854R$, or the following data exist:
\begin{enumerate}
\item A number field $L$ with $[L:\Q]\le576$, and a scalar
character of order $m$ satisfying
\[
 m\le2\prod_{\substack{p\mid R\\p\text{ odd}}}(p-1).
\]
\item A normal field $K/\Q$ of degree
$d'\le2m[L:\Q]$ and a primitive weight-two newform of level $N$,
with nebentypus $\varepsilon$, such that
\[
 N\le CR^2,\qquad \rad(N)\mid R,
 \qquad\varepsilon^2|_{G_K}=1.
\]
\item For $H=\ker\varepsilon\subset(\Z/N\Z)^\times$, the
abelian surface $A_K$ is an isogeny factor over $K$ of
$J_H=\Jac(X_1(N)/H)_K$.
\end{enumerate}
\end{definition}

The first alternative is supplied by the complex multiplication
case. In the second alternative, the square of the nebentypus is
trivial over $K$ by the determinant identity for the compatible
representations. The newform quotient factors through $J_H$ because
every element of $H$ fixes all conjugates of its differentials.
This is an assertion about the order of a character; it does not
assert a bounded degree for the newform's coefficient field.

\begin{theorem}\label{thm:contact-height}
Assume the modular realization input. Define
\[
 H_{18}=2^{110}210^{1800}B_*^{18}+128.
\]
Then
\[
 h(a)\le H_{18}M_5(a)^{18}
 \quad\text{for every }a\in\Q\setminus\mathcal D,
 \qquad H_{18}<2^{1984521}.
\]
\end{theorem}

\section{The degree of the modular curve}

\begin{lemma}\label{lem:degree}
In the non-CM alternative of Definition~\ref{def:input},
\[
 d'\le2304\varphi(R),\qquad
 t:=\ord(\varepsilon)\le4608\varphi(R),
\]
and
\begin{equation}\label{eq:degree}
 d_H:=\deg(X_H\longrightarrow X(1))
 =Nt\prod_{p\mid N}(1+p^{-1})
 <4608CR^3.
\end{equation}
The genus $g$ of $X_H$ satisfies $2\le g\le d_H$.
\end{lemma}
\begin{proof}
The integer $R$ is squarefree and even. Thus
\[
 m\le2\prod_{\substack{p\mid R\\p\text{ odd}}}(p-1)
      =2\varphi(R),
 \qquad
 d'\le1152m\le2304\varphi(R).
\]
Since $\varepsilon^2$ is trivial on $G_K$, the order of its image
divides $d'$. Therefore $t\le2d'\le4608\varphi(R)$.

The congruence subgroup defining $X_H$ is
\[
 \Gamma_H=\left\{
 \begin{pmatrix}a&b\\c&d\end{pmatrix}\in\Gamma_0(N):
 d\bmod N\in H\right\}.
\]
A weight-two nebentypus is even, so $-I\in\Gamma_H$. Consequently
the degree of the coarse map to $X(1)$ equals its
$\operatorname{SL}_2(\Z)$ index, with no extra factor of two.
The standard index formula gives the equality in \eqref{eq:degree};
see \cite[Section~4, equations~(4.1)--(4.2)]{vonKaenel}.

Since $\rad(N)\mid R$, the exact totient identity gives
\[
 \begin{split}
 d_H
 &\le4608CR^2\varphi(R)\prod_{p\mid R}(1+p^{-1})\\
 &=4608CR^3\prod_{p\mid R}(1-p^{-2})
 <4608CR^3.
 \end{split}
\]
The map is branched only over the cusp and the two elliptic
points. It is therefore a Belyi map after moving those three
points to $0,1,\infty$. The usual genus formula yields
$g\le1+d_H/12\le d_H$. Finally $A$ is a surface isogeny
factor of $J_H$, so $g\ge2$.
\end{proof}

\section{Faltings height and the isogeny comparison}

Set
\[
 T=2^{14}CR^3.
\]
Lemma~\ref{lem:degree} gives $g<T$ and $d_H<T$, and $T\ge16$.
Javanpeykar's Belyi-curve estimate
\cite[Theorem~1.1.1 and Section~2.3]{Javanpeykar} implies
\begin{equation}\label{eq:jacobian}
 \hF(J_H)\le13000000\,d_H^5g<2^{24}T^6.
\end{equation}
This estimate applies to every Belyi curve. It requires neither
squarefree level nor trivial nebentypus.

\begin{lemma}\label{lem:isogeny}
The abelian surface in Definition~\ref{def:input} satisfies
\[
 \hF(A)<2^{25}T^6=2^{109}210^{1800}R^{18}.
\]
\end{lemma}
\begin{proof}
In the non-CM case, Poincar\'e reducibility gives an abelian variety
$Z/K$ such that $J_H$ is $K$-isogenous to $A\times Z$.
Apply the isogeny theorem of Gaudron--R\'emond
\cite[Theorem~1.9(1), pp.~7--8]{GaudronRemond} with source $J_H$
and target $A\times Z$. Half-logarithmic isogeny-height variation
and additivity of Faltings height give
\[
 \hF(A)+\hF(Z)\le\hF(J_H)+V,
\]
where
\[
 V=2g^2\left(8g^2\log(7g)+\log d'
       +\log\max(1,\log d',\hF(J_H))\right).
\]
The Bost lower bound, in the same normalization, gives
$\hF(Z)>-3g/2$; see \cite[equation~(2.2)]{vonKaenel}.
Write $u=\log T$. The inequalities $2304<e^8$, $7<e^2$,
and $2^{24}<e^{17}$ imply
\[
 \log d'<8+u,\qquad \log(7g)<2+u,
\]
and, using \eqref{eq:jacobian},
\[
 \log\max(1,\log d',\hF(J_H))<17+6u.
\]
Hence
\begin{equation}\label{eq:V}
 V\le16T^4(2+u)+2T^2(25+7u)<T^6.
\end{equation}
For completeness, $\log T\le T/4$ for $T\ge16$. Divide the
first bound in \eqref{eq:V} by $T^6$ to obtain
\[
 \frac{V}{T^6}
 \le\frac4T+\frac{32}{T^2}+\frac{7}{2T^3}+\frac{50}{T^4}
 \le\frac4{16}+\frac{32}{16^2}
           +\frac7{2\cdot16^3}+\frac{50}{16^4}<1.
\]
It follows that
\[
 \hF(A)<(2^{24}+1)T^6+\tfrac32T<2^{25}T^6.
\]
The CM estimate $\hF(A)<854R$ is also smaller than this bound.
\end{proof}

\begin{proof}[Proof of Theorem~\ref{thm:contact-height}]
Combine \eqref{eq:geometric} with Lemma~\ref{lem:isogeny} and
use $R=B_*M_5(a)$ and $M_5(a)\ge1$. This gives
\[
 h(a)\le
 \left(2^{110}210^{1800}B_*^{18}+128\right)M_5(a)^{18}.
\]
The exact integer computation described below verifies
$B_*<2^{109445}$. Since $210<2^8$ and
$2^{110}210^{1800}B_*^{18}>128$,
\[
 H_{18}
 <2^{111}210^{1800}B_*^{18}
 <2^{111+1800\cdot8+18\cdot109445}
 =2^{1984521}.
\]
\end{proof}

\section{Disjoint sets of contacts}

\begin{corollary}\label{cor:amplification}
Fix a positive integer $r$. For
\[
 \Theta_r=\bigcup_{j=0}^{r-1}
 \left(3j+\left\{0,-\frac12,-1,
                       \frac1{2399},-\frac{32}{407}\right\}\right)
\]
there is a fixed finite exceptional set of primes $\Sigma_r$
and a constant $C_r>0$ such that
\[
 h(t)\le C_r M_r(t)^{18/r}
 \qquad(t\in\Q\setminus\Theta_r).
\]
Here $M_r(t)$ is the product of primes outside $\Sigma_r$ at
which $t-\theta$ has positive odd valuation for some
$\theta\in\Theta_r$.
\end{corollary}
\begin{proof}
Apply Theorem~\ref{thm:contact-height} to
$\phi_j(t)=2+1/(t-3j)$ for $0\le j<r$.
Each matrix has determinant $-1$, and therefore
$h(t)\le h(\phi_j(t))+O_r(1)$.
Remove the primes in $S_*$, those dividing the coefficients
needed to clear contact denominators, and those dividing
pairwise determinants of the $5r$ primitive contact forms.
Outside this fixed finite set, different contact forms cannot
vanish at the same residue point. The radicals for different
blocks thus have disjoint prime support. Primes removed from
an individual radical contribute at most a fixed factor.
Multiplication of the $r$ resulting inequalities gives the claim.
Explicit contact forms and their determinants are recorded in
\cite[Sections~1--2]{Amplification}.
\end{proof}

\section{Exact numerical verification}

The script
\begin{center}\small
\nolinkurl{computation/contact_heights/check_contact_height18.py}
\end{center}
reconstructs the prime product $B_*$, the integer $H_{18}$,
and their binary lengths. It also verifies the rational bound
in \eqref{eq:V}, the numerical exponent identity, and the three
exponential inequalities used in the isogeny estimate by rational
partial sums of the exponential series. Its certificate is
\nolinkurl{contact-height18-certificate.json} in the same directory.
The symbolic totient identity and the arithmetic implications
are proved above; they are not inferred from a bounded search.

\begin{thebibliography}{9}
\bibitem{Javanpeykar}
A.~Javanpeykar, with an appendix by P.~Bruin,
\emph{Polynomial bounds for Arakelov invariants of Belyi curves},
Algebra \& Number Theory \textbf{8} (2014), 89--140,
\url{https://doi.org/10.2140/ant.2014.8.89}.

\bibitem{GaudronRemond}
\'E.~Gaudron and G.~R\'emond,
\emph{Nouveaux th\'eor\`emes d'isog\'enie},
M\'emoires de la Soci\'et\'e Math\'ematique de France
\textbf{176} (2023), Theorem~1.9(1),
\url{https://doi.org/10.24033/msmf.484}.

\bibitem{vonKaenel}
R.~von K\"anel, \emph{The effective Shafarevich conjecture for abelian
varieties of $\mathrm{GL}_2$-type}, Forum of Mathematics, Sigma
\textbf{9} (2021), e39,
\url{https://doi.org/10.1017/fms.2021.29}.

\bibitem{Glazer}
E.~Glazer and contributors,
\emph{Diophantine definition of the integers in the rationals},
\texttt{manuscripts/quartic\_v9.tex}, commit
\texttt{14e8f601fe0d9173d4a5f084aedc3e07066a7688},
\url{https://github.com/elliotglazer/diophantine-z-q}.

\bibitem{Amplification}
J.~Brox and contributors,
\emph{Height bounds from disjoint sets of rational contacts},
\texttt{contact\_height\_amplification\_20261109\_v13.tex},
\url{https://github.com/JoseBrox/diophantine-z-q_improvements}.
\end{thebibliography}
\end{document}
